% A.40: independently audited and registered by the parent agent.
\subsection{All genus-one covers}

\begin{statement}{A.40}{Proved}{Every genus-one cover has exactly cusp fixed groups and cusp normalizers}
Let \(Y_H\to U\) be any finite connected geometric cover
whose smooth compactification has genus one. After a finite
field extension and a choice of one cusp as origin, its
compactification map is an isogeny
\[
                \varphi:E'\longrightarrow E,
       \qquad S=\ker\varphi,
       \qquad Y_H=E'\setminus S.
\]
For every prime \(\ell\), every original packet, and every
permitted open subgroup \(\Delta_0\) preserving the cover
and its selected cusp lift \(P\),
\[
       \left(\pi_1(Y_{H,\bar K})^{(\ell)}\right)^{s_b(\Delta_0)}
                              =I_{P,\ell}.             \tag{A.40.1}
\]
Every nontrivial closed procyclic subgroup of this pro-\(\ell\)
group whose arithmetic normalizer has open Galois image is
contained in a conjugate of a cusp inertia group.
Thus Task A.5 holds for every genus-one compactification,
with no restriction on the isogeny degree.

The fixed-group proof includes the following more general
geometric criterion. Suppose a smooth proper curve of positive
genus has potentially good reduction at a finite place of
residue characteristic \(p\), and two distinct cusps, one of
them the selected cusp, have the same specialization in its
good model. For \(\ell\ne p\), the leading weighted Lie
classes of its actual derived-fixed subgroup lie in the free
Lie algebra generated by the boundary classes and the total
proper symplectic class. Consequently that fixed subgroup is
the selected inertia, by A.38. Other cusps may specialize at
different points; their locations do not affect the assertion.
\end{statement}

\begin{proof}
\paragraph{A single positive-genus component and its marked trees.}
First prove the geometric criterion. Enlarge the finite number
field \(k\) so all cusps and a pointed semistable model at the
chosen place are defined, the proper curve has good reduction,
and \(\Gamma=\overline{[G_k,G_k]}\) lies in the specified
\(\Delta_0\), when one is prescribed. Let \(P\ne Q\) be
the two cusps with equal specialization.

Separating the marked sections in the smooth proper model,
and then stabilizing, produces one positive-genus component
\(C\), with rational marked trees attached at its smooth
points. Every finite edge has a strictly positive rational
length. The cusps \(P,Q\) lie in the same attached tree.
This follows directly by blowing up the finitely many common
specializations of sections: a blowup introduces a rational
exceptional component and preserves the tree structure;
contracting an unstable rational component adds the lengths
of its adjacent edges. In particular the paths from \(P\)
to \(C\) and to \(Q\) separate strictly before reaching
\(C\).

Root this metric tree at the component carrying \(P\).
Write \(e\) for its finite edges, \(m_e>0\) for their
lengths, and orient each edge away from the root. A cusp
already on \(C\) is simply assigned to that vertex. Let
\(E_e\), \(Q_e\), and \(j_e\) be the indicators that
\(C\), \(Q\), and cusp \(j\), respectively, lie on the
side of \(e\) away from the root.

Use handle generators \(a_i,b_i\), \(1\le i\le g\),
and cusp generators \(t_j\), of weights one and two
respectively. Set
\[
                  \Omega=\sum_{i=1}^g[a_i,b_i].
\]
The first peripheral relation is
\(\Omega+\sum_jt_j=0\). Eliminate \(t_Q\), keeping
the selected \(t_P\) as an independent generator. The
weighted graded Lie algebra is the free Lie algebra on
the handles and these remaining cusp generators.

For an ordinary base point in the selected cusp disc,
the actual based edge-twist formula of
\cite[Definition 3.1.5 and Theorem 3.1.6]{fghBettsDogra}
gives the first weighted monodromy derivation. Put
\[
 \delta_e=E_e\Omega+\sum_{j:\,j_e=1}t_j,
 \qquad
 W=\sum_e m_eE_e\delta_e,
 \qquad W_j=\sum_e m_ej_e\delta_e.                    \tag{A.40.2}
\]
With a common orientation convention it has
\[
         Nv=[W,v]\quad\text{for every handle }v,
             \qquad Nt_j=[W_j,t_j].                 \tag{A.40.3}
\]
Changing the common orientation negates the entire
derivation and has no effect on its kernel. All
edge boundaries have weight two. The comparison with
the chosen tangent takes place in its actual cusp
disc. As in A.36 and A.39, an arbitrary packet changes
the genuine first derivation only by
\(\lambda\operatorname{ad}(t_P)\),
\(\lambda\in\mathbf Q_\ell\). Thus replace
\(W,W_j\) by \(W+\lambda t_P,W_j+\lambda t_P\).
This change leaves every coefficient of \(\Omega\)
below unchanged. There is no weight-one discrepancy:
such an inner discrepancy would centralize \(t_P\)
in the free Lie algebra, whose centralizer there
contains no weight-one element.

After elimination of \(Q\), the coefficient of
\(\Omega\) in \(\delta_e\) is \(E_e-Q_e\).
Consequently the coefficient in \(W\), denoted \(r\),
and that in \(W_j\), denoted \(c_j\), are
\[
 r=\sum_e m_eE_e(1-Q_e),\qquad
 c_j=\sum_e m_ej_e(E_e-Q_e).                         \tag{A.40.4}
\]
The path from the branching point of \(P,Q,C\) to
\(C\) has positive length, and its edges are precisely
those counted by \(r\). Hence \(r>0\). Some \(c_j\)
can be negative or equal to \(r\); no sign or strict
inequality for them will be used. In the special case
where all cusps lie in one attached tree, one also has
\[
 r-c_j=\sum_e m_e\bigl(E_e(1-Q_e)(1-j_e)
                   +(1-E_e)Q_ej_e\bigr)>0,          \tag{A.40.5}
\]
since the last edge incident to \(C\) contributes.

An actual local Frobenius--tame-inertia commutator in
\(\Gamma\) has first derivation a nonzero common
multiple of (A.40.3) with this packet correction.
The multiplier is \(q_{\rm res}-1\ne0\); after the
semistable extension, the prime-to-\(p\) monodromy is
the actual edge action just used. Thus it suffices to
bound the kernel of this genuine rational derivation.

\paragraph{The unbounded prefix calculation.}
Let \(A\) be the free associative algebra on the handles
and retained cusp letters, and put
\[
               B=\mathbf Q_\ell\langle\Omega,t_j:
                                                j\ne Q\rangle.
\]
It is a free algebra on the displayed generators.
For completeness, the leading word of \(\Omega\)
is \(a_1b_1\), with coefficient one and no proper
self-overlap, after choosing an appropriate order.
The handle algebra has the triangular basis
\[
 \{\Omega^n:n\ge0\}\ \sqcup
 \{\Omega^u w\Omega^v:u,v\ge0,\quad
         w\ne\varnothing\text{ begins and ends with no }a_1b_1\}.
                                                               \tag{A.40.6}
\]
Indeed, strip initial and terminal powers of
\(a_1b_1\) from each handle word; the indicated
leading words are distinct and exhaust all words.
Adjoining the cusp letters by free product gives a
normal basis whose words outside \(B\) have a first
nonboundary handle core with an arbitrary \(B\)-prefix.

Let \(D_B\) denote the restriction of the corrected
derivation to \(B\). On every handle word the
derivation is the commutator with the same \(W\),
including its packet correction. On the prefix of
the first core, its portion increasing boundary-word
length is therefore
\[
                            K=D_B+R_W.              \tag{A.40.7}
\]
All terms acting after that core preserve its prefix
length. They cannot remove the core: left or right
multiplication by \(\Omega\) remains in its displayed
bimodule, and a cusp letter merely adds a free-product
separator. Thus injectivity of \(K\) excludes all
elements outside \(B\) from \(\ker D\).

Filter \(B\) by the number of occurrences of
\(\Omega\). Since
\(D_B\Omega=[W,\Omega]\), this term does not raise
that count. The part of \(K\) which raises it by one
is
\[
  K^+=D_c+rR_\Omega,
  \qquad D_c\Omega=0,
  \qquad D_ct_j=c_j[\Omega,t_j].                    \tag{A.40.8}
\]
All remaining terms, including every packet term,
preserve the count.

Fix the ordered cusp word \(i_1\cdots i_n\). Its
span of insertions
\[
          \Omega^{a_0}t_{i_1}\Omega^{a_1}
                         \cdots t_{i_n}\Omega^{a_n}
\]
identifies linearly with the polynomial ring
\(\mathbf Q_\ell[Z_0,\ldots,Z_n]\), by sending
the insertion to \(Z_0^{a_0}\cdots Z_n^{a_n}\).
Under this identification \(K^+\) is multiplication
by
\[
 c_{i_1}Z_0+
       \sum_{h=1}^{n-1}(c_{i_{h+1}}-c_{i_h})Z_h
                       +(r-c_{i_n})Z_n.             \tag{A.40.9}
\]
Its coefficients sum to \(r\ne0\), so this
polynomial is nonzero, regardless of the signs of
the \(c_j\). For \(n=0\), the multiplier is simply
\(rZ_0\). A polynomial ring over \(\mathbf Q_\ell\)
is a domain, so \(K^+\) is injective on every one
of these mutually independent cusp-word summands.
It follows that \(K\) is injective by taking the
largest \(\Omega\)-count in a nonzero polynomial.
Taking a largest first-core prefix length now proves
\[
                \ker D\subset B\cap\mathbb L
                    =\mathbb L(t_j:j\ne Q,\Omega).  \tag{A.40.10}
\]
The last equality is the primitive-element identity
for the tensor Hopf algebra on these freely generated
primitive elements. The argument works in every
weight; it is not a finite-degree dimension count.

Every leading weighted class of the unipotent closure
of the actual \(\Gamma\)-fixed subgroup is killed
by \(D\), hence lies in (A.40.10). Each retained
cusp class has character \(\chi_\ell\), and so does
\(\Omega\), by the proper symplectic similitude.
All its generators have weight two. Thus the closure
is Tate-graded in the precise sense of A.38, which
gives the selected inertia as the actual fixed group.
This proves the geometric criterion.

\paragraph{A collision for every prime dividing an isogeny degree.}
Return to \(\varphi:E'\to E\), and let
\(N=\deg\varphi>1\). We claim that for every
prime \(p\mid N\), some place above \(p\), after
a finite extension, has a nonzero element of
\(S[p^\infty]\) specializing to zero in the good
model of \(E'\).

Both curves acquire good reduction
by \cite[Corollary VII.7.2]{fghSilvermanAEC}, and have the
same ordinary or supersingular reduction type after
such an extension: their isogeny extends between
the good models. The degree of this special isogeny
equals its generic degree, since the degree of the
pullback of an ample degree-one line bundle is
constant in a smooth proper family. At \(p=2\),
\(E\) is potentially supersingular, as already used
in A.39. At odd \(p\equiv3\pmod4\), it is
supersingular by
\cite[Theorem V.4.1(a), Example V.4.5]{fghSilvermanAEC}:
the expansion of \((x^3-x)^{(p-1)/2}\) has no
term \(x^{p-1}\). In both cases the geometric
special fibre has no nonzero \(p\)-power torsion,
so every point of \(S[p^\infty]\) specializes to
zero, proving the claim.

Suppose now \(p\equiv1\pmod4\). Write
\(p=\pi\bar\pi\) in \(\mathbf Z[i]\), and let
\(p^n\) be the exact \(p\)-part of \(N\).
The reductions are ordinary. At a place of \(K\)
dividing \(\pi\), the isogeny \([\pi]\) of \(E\)
has degree \(p\) and zero special differential, so
its reduction is purely inseparable. The isogeny
\([\bar\pi]\) has invertible differential and
separable reduction. Therefore, among
\(E[p^n]\), the points specializing to zero are
exactly \(E[\pi^n]\). At the conjugate prime they
are exactly \(E[\bar\pi^n]\). This can also be
seen by decomposing
\(E[p^n]=E[\pi^n]\oplus E[\bar\pi^n]\): the
first summand has purely inseparable reduction and
the second is finite etale at the selected place.

Assume for contradiction that \(S[p^\infty]\)
specializes injectively at every place above \(p\).
Prime-to-\(p\) torsion always specializes injectively
in good reduction. Thus all \(N\) kernel points
remain distinct in the special kernel, and
\(\widetilde\varphi\) is separable of degree \(N\).
The dual isogeny satisfies
\(\widehat\varphi\varphi=[N]\). In ordinary
characteristic \(p\), \([N]\) has inseparable
degree \(p^n\); hence \(\widehat\varphi\) has
inseparable degree \(p^n\) and separable degree
\(N/p^n\). In particular the special kernel of
\(\widehat\varphi\) has no nonzero geometric
\(p\)-primary point. Its generic \(p\)-primary
kernel \(T\subset E[p^n]\), which has order
\(p^n\), therefore specializes to zero at both
CM primes. It follows that
\[
                  T\subset E[\pi^n]\cap E[\bar\pi^n]=0,
\]
because the two ideals are coprime. This contradicts
\(|T|=p^n>1\) and proves the collision claim.
The dual-isogeny and inseparable-degree facts here
are \cite[Theorem III.6.1, Corollary III.6.4,
Theorem V.3.1]{fghSilvermanAEC}; prime-to-\(p\)
specialization is \cite[Proposition VII.3.1(b)]{fghSilvermanAEC}.

Translation by any selected cusp \(P\in S\)
extends to the good model. Thus a nonzero kernel
point specializing to zero gives a distinct cusp
\(Q\in S\) specializing to the same point as
\(P\). The collision claim applies to every
selected cusp, not just the chosen origin.

\paragraph{All coefficient primes and the normalizers.}
Fix \(\ell\). If \(N\) has a prime divisor
\(p\ne\ell\), choose its collision place and
apply the geometric criterion above. If \(N\)
is an \(\ell\)-power, the geometric isogeny
cover has Galois group \(S\), an \(\ell\)-group,
so A.22 gives (A.40.1) directly. Degree one is
the original curve and is covered by R1, with
the same open-subgroup strengthening used in A.22.
This proves (A.40.1) for every isogeny degree,
every coefficient prime, and every stated packet.

For the normalizer assertion, first enlarge the
finite field of definition so the isogeny, the
cusp origin, and all required paths are defined.
The map \(Y_H\to U\) induces a continuous map
of geometrically pro-\(\ell\) arithmetic groups
over the identity of this open Galois group. A
closed procyclic subgroup with open normalizer
therefore has normalized image with open arithmetic
image whenever that image is nontrivial.

If the source subgroup were noncuspidal, A.11
would give it a nonzero class in the proper
\(\mathbf Q_\ell\)-homology of \(E'\). The
isogeny induces an isomorphism
\(V_\ell E'\to V_\ell E\), even when
\(\ell\mid N\). Its image subgroup is thus
nontrivial and has nonzero proper homology.
By R2, independently verified in A.11--A.12,
the target subgroup is cuspidal, which gives
zero proper homology, a contradiction. Restriction
to a finite field leaves the original normalizer
image open, so this proves the stated classification.

Finally every finite map between the genus-one
compactifications is unramified by Riemann--Hurwitz.
Choosing a point above \(O\) as origin makes it
an isogeny, and its removed set is exactly its
kernel. These cover all genus-one cases in the
statement.

The program \texttt{code/a\_isogeny\_fixed.py}
checks the rooted-tree coefficient identity on
all labelled trees of small orders, the exact
slot polynomial formula, and the prefix injectivity
on an explicit branched marked tree, including
negative coefficients and arbitrary sample packet
scalars. The all-degree and all-tree assertions
are proved above, independently of those checks.
\end{proof}
